%======================================================================
\chapter{No arc, and hereditary indecomposability}\label{ch:noarc}
%======================================================================

This chapter proves that a homogeneous indecomposable plane continuum contains no arc (Bing 1960 for the
indecomposable case; F.~B.~Jones, \emph{Use of a new technique in homogeneous continua}, 1975), and deduces from
it and Chapter~\ref{ch:decomp} Hagopian's theorem that such a continuum is hereditarily indecomposable.

Jones's paper proves the no-arc theorem from six properties (a)--(f) of a homogeneous indecomposable plane
continuum containing an arc. The formal proof keeps his overall idea (study the arc components, which are lines,
and use the Effros theorem to transport pieces of a line along itself) but replaces two of his steps, whose
printed proofs are defective, and replaces his final planar argument by a different one.

\begin{itemize}
\item Jones's property (f) (``no folding in a thin neighbourhood'') is proved from (a)--(c) and the Effros theorem
alone. That argument cannot be right: the dyadic solenoid satisfies everything it uses, and (f) fails there. The
formal proof replaces (f) by two properties, (R) and (O) below, which do hold in that generality.
\item Jones's argument that every point of a line is a limit of the rest of the line has a gap (the Effros
homeomorphism might reverse the line); it is replaced by a short recurrence argument.
\item The final planar contradiction is a ``trapping'' argument in a collar of an arc, in the spirit of the
Poincar\'e--Bendixson theorem, using the Jordan curve theorem.
\end{itemize}

\section{The setting}

\begin{definition}\label{def:jones-setting}
A \emph{Jones setting} is a compact metric space $Y$ that is homogeneous, hereditarily unicoherent and atriodic and
contains an arc.
\end{definition}

\formal{\lean{HPC.Hagopian.JonesSetting}, \lean{HPC.Hagopian.jonesSetting_of_plane}.}

By Theorem~\ref{thm:hagopian-lemma1}, a homogeneous indecomposable plane continuum that contains an arc is a Jones
setting. Sections~\ref{sec:lines}--\ref{sec:recurrence} work in an arbitrary Jones setting $Y$; planarity enters
only in Section~\ref{sec:trap}. (The dyadic solenoid is a Jones setting, so these sections are not vacuous.)

\section{Arcs and lines}\label{sec:lines}

\begin{proposition}[Jones (c)]\label{prop:jones-c}
Let $A,B$ be arcs in $Y$ that meet. Then $A\cup B$ is an arc. If $A$ and $B$ have a common end point and meet in another
point, one of them contains the other.
\end{proposition}

\begin{proof}
$A\cap B$ is a continuum (hereditary unicoherence), hence a point or a sub-arc of each. At each end point of $A\cap B$
count the arms of $A\cup B$ there: the sub-arc $A\cap B$ (if nondegenerate), a tail of $A$, a tail of $B$. Three arms that
pairwise meet only at that point form a triod (\vthree{Lemma~4.5}), which is excluded. So there is at most one tail at
each end (at most two in all if $A\cap B$ is a point), and $A\cup B$ is an arc. For the second statement, $A\cap B$ is a sub-arc
from the common end point $a$; its far end is an end point of $A$ or of $B$, since otherwise there would be three arms
there.
\end{proof}

\formal{\lean{HPC.Hagopian.JonesSetting.inter_isPreconnected}, \lean{HPC.Hagopian.JonesSetting.no_three_arms},
\lean{HPC.Hagopian.JonesSetting.union_arc}, \lean{HPC.Hagopian.JonesSetting.nested}.}

In particular two arcs with the same end points are equal. By homogeneity every point is an interior point of an arc
(map an interior point of a given arc to it).

\formal{\lean{HPC.Hagopian.JonesSetting.arc_unique}, \lean{HPC.Hagopian.JonesSetting.exists_arc_through}.}

The \emph{arc component} of $x$ is the union of $\set x$ and all arcs containing $x$. A \emph{line} is a continuous
injection $\ell\colon\R\to Y$ whose image is an arc component.

\formal{\lean{HPC.Hagopian.arcComp}, \lean{HPC.Hagopian.IsLine}.}

\begin{theorem}\label{thm:lines}
\leavevmode
\begin{enumerate}[label=(\alph*)]
\item Every arc component of $Y$ is the image of a line.
\item If $\ell$ is a line and $s<t$, then $\ell[s,t]$ is an arc, and every arc between two points of $\ell$ is of this form.
\item If $\ell,\ell'$ are lines and $h$ is a homeomorphism of $Y$ with $h(\ell(0))\in\ell'(\R)$, then $h\circ\ell=\ell'\circ\psi$ for a
continuous strictly monotone bijection $\psi$ of $\R$.
\end{enumerate}
\end{theorem}

\begin{proof}
(a) Fix an arc $[a,c]$ with $x$ in its interior. The arcs from $x$ containing $[x,c]$ form a chain under inclusion
(Proposition~\ref{prop:jones-c}). There is no largest one: if $[x,z]$ were largest, an arc with $z$ in its interior would
extend it beyond $z$. The union $R^+$ of the chain is separable; a countable dense subset is covered by countably many
arcs of the chain, and these are cofinal (if all lay in one arc $[x,y]$, $R^+$ would lie in the closed set $[x,y]$, which
would be largest). Parametrize $[0,\infty)$ piecewise along an increasing cofinal sequence. Do the same for $R^-$ through
$a$. $R^+\cap R^-=\set x$, since an arc from $x$ containing $[x,a]$ and $[x,c]$ would make them nested, and $R^+\cup R^-$ is
the arc component by Proposition~\ref{prop:jones-c}. Gluing the two parametrizations gives a continuous injection of $\R$.

(b) The first statement holds for any continuous injection; the second because two arcs with the same end points are
equal.

(c) $h$ maps arc components onto arc components, so it maps $\ell(\R)$ onto $\ell'(\R)$; $\psi=\ell'^{-1}\circ h\circ\ell$ is a bijection of $\R$
that maps compact intervals onto compact intervals by (b), hence is continuous and strictly monotone.
\end{proof}

\formal{\lean{HPC.Hagopian.JonesSetting.rayArcs_total}, \lean{HPC.Hagopian.JonesSetting.extend},
\lean{HPC.Hagopian.JonesSetting.exists_ray_seq}, \lean{HPC.Hagopian.JonesSetting.exists_line},
\lean{HPC.Hagopian.JonesSetting.line_arc}, \lean{HPC.Hagopian.JonesSetting.arc_in_line},
\lean{HPC.Hagopian.JonesSetting.homeo_on_line}.}

\section{No small reversal}\label{sec:reversal}

\begin{lemma}[(R0)]\label{lem:R0}
Let $A=[a,p]$ and $C=[p,c]$ be arcs and $\rho=\min(\dist(a,C),\dist(c,A))$. If $h$ is a homeomorphism with $h(p)=p$ that moves every
point by less than $\rho$, then $h(A)\cap C=\set p$.
\end{lemma}

\begin{proof}
Suppose $h(A)$ meets $C\setminus\set p$. The arcs $h(A)$ and $C$ start at $p$ and meet elsewhere, so they are nested
(Proposition~\ref{prop:jones-c}). If $C\subseteq h(A)$, then $c=h(y)$ with $y\in A$ and $d(c,y)<\rho\le\dist(c,A)$; if $h(A)\subseteq C$,
then $h(a)\in C$ and $d(a,h(a))<\rho\le\dist(a,C)$. Both are impossible.
\end{proof}

\formal{\lean{HPC.Hagopian.JonesSetting.no_reversal_arc}.}

\begin{theorem}[(R) and (O)]\label{thm:RO}
There is $\varepsilon_0>0$ with the following properties.
\begin{enumerate}[label=(\alph*)]
\item \emph{(R)} If $h$ is an $\varepsilon_0$-homeomorphism fixing $\ell(s)$ for a line $\ell$, then $h$ maps $\ell(t)$, $t>s$, to points
$\ell(t')$ with $t'>s$.
\item \emph{(O)} If $h$ is an $\varepsilon_0$-homeomorphism with $h(\ell(0))\in\ell(\R)$, then $h\circ\ell=\ell\circ\psi$ with $\psi$ strictly
increasing, continuous and onto.
\end{enumerate}
\end{theorem}

\begin{proof}
(a) For each $q\in Y$ take an arc with $q$ in its interior and its $\rho_q$ as in Lemma~\ref{lem:R0}, and let $\delta_q$ be the
Effros bound for $\rho_q/3$. If $d(b,q)<\delta_q$, let $g$ be a $\rho_q/3$-homeomorphism with $g(b)=q$. If $h$ fixes $b$, moves points
less than $\rho_q/3$ and reverses the line through $b$, then $ghg^{-1}$ fixes $q$, moves points less than $\rho_q$, and reverses
the line through $q$; by Theorem~\ref{thm:lines}(c) reversing the line at a fixed point means reversing every arc there,
which contradicts Lemma~\ref{lem:R0}. Cover $Y$ by finitely many balls $B(q,\delta_q)$ and take $\varepsilon_0=\min\rho_q/3$.

(b) By Theorem~\ref{thm:lines}(c), $h\circ\ell=\ell\circ\psi$ with $\psi$ a continuous monotone bijection. If $\psi$ decreased, $\psi(s)-s$ would be
continuous, strictly decreasing and onto $\R$, so $\psi(s^*)=s^*$ for some $s^*$, and $h$ would fix $\ell(s^*)$ and reverse $\ell$
there, against (a).
\end{proof}

\formal{\lean{HPC.Hagopian.JonesSetting.local_no_reversal}, \lean{HPC.Hagopian.JonesSetting.exists_no_reversal},
\lean{HPC.Hagopian.JonesSetting.orientation}.}

\begin{remark}
Jones's (f) asserts, for a point $p$ and an open set $D$ with $p\notin\cl D$, an $\varepsilon$ such that no arc $abc$ with $a,c\in D$ has
the $\varepsilon$-neighbourhood of $ab$ containing $p$ and $bc$. His proof uses only the Effros theorem and (c). In the dyadic
solenoid, where all of these hold, (f) fails: a long initial piece $\lambda[0,S]$ of a line is $\varepsilon/2$-dense, and an arc $abc$
along the line beyond it is a counterexample. Theorem~\ref{thm:RO} is what the argument actually needs.
\end{remark}

\section{Recurrence}\label{sec:recurrence}

For a line $\ell$, its \emph{forward limit set} is $\omega(\ell)=\set{z: z=\lim\ell(u_n)\text{ for some }u_n\to\infty}$.

\formal{\lean{HPC.Hagopian.omegaLim}, \lean{HPC.Hagopian.mem_omegaLim_iff}.}

\begin{theorem}\label{thm:recurrent}
\leavevmode
\begin{enumerate}[label=(\alph*)]
\item $\omega(\ell)$ contains the arc component of each of its points.
\item Every line lies in its own forward limit set: $\ell(\R)\subseteq\omega(\ell)$. The same holds backwards.
\end{enumerate}
\end{theorem}

\begin{proof}
(a) Let $z=\lim\ell(u_n)$, $u_n\to\infty$, and $\lambda$ the line through $z$ with $\lambda(0)=z$. The Effros theorem gives
$\varepsilon_n$-homeomorphisms $h_n$, $\varepsilon_n\to0$, with $h_n(z)=\ell(u_n)$; by Theorem~\ref{thm:lines}(c), $h_n(\lambda(s))=\ell(\psi_n(s))$
with $\psi_n(0)=u_n$. Fix $s$. If $\psi_n(s)\le C$ along a subsequence, then for every $L$, eventually $\ell[C,C+L]\subseteq h_n(\lambda[0,s])$ (the
parameter interval between $\psi_n(s)$ and $u_n$ eventually covers $[C,C+L]$), which lies within $\varepsilon_n$ of $\lambda[0,s]$. So
$\ell[C,\infty)\subseteq\lambda[0,s]$, a compact arc, which is impossible (the arc components of $\ell$ and $\lambda$ are equal or disjoint, and
$\ell$ is injective). Hence $\psi_n(s)\to\infty$ and $\lambda(s)\in\omega(\ell)$.

(b) $K=\omega(\ell)$ is nonempty and compact. Suppose $\ell(0)\notin K$ and let $\varepsilon=\frac12\dist(\ell(0),K)$. Take $p\in K$, $\delta$ the Effros
bound for $\varepsilon$, $t$ with $d(\ell(t),p)<\delta$, and an $\varepsilon$-homeomorphism $g$ with $g(\ell(t))=p$. $g$ maps the line $\ell$ onto the
line through $p$, which lies in $K$ by (a); so $g(\ell(0))\in K$, while $d(\ell(0),g(\ell(0)))<\varepsilon<\dist(\ell(0),K)$. Apply this to the shifted
lines $s\mapsto\ell(s+c)$ and the reversed line $s\mapsto\ell(-s)$.
\end{proof}

\formal{\lean{HPC.Hagopian.JonesSetting.omegaLim_saturated}, \lean{HPC.Hagopian.JonesSetting.zero_mem_omegaLim},
\lean{HPC.Hagopian.JonesSetting.line_recurrent}, \lean{HPC.Hagopian.JonesSetting.jones75_d}.}

\section{Passes through a box}\label{sec:passes}

From now on $M$ is a homogeneous indecomposable plane continuum containing an arc, $\ell$ is a line in $M$, and
$\gamma=\ell|_{[0,1]}$. By the collar theorem (\vfive{Theorem~9.19}) there is a \emph{collar} $F$ of $\gamma$: an embedding of
$[0,1]\times[-1,1]$ into the plane with $F(t,0)=\gamma(t)$ and $F((0,1)\times(-1,1))$ open. On $F([0,1]\times[-1,1])$ we use the collar coordinates
$(X,Y)=F^{-1}$. The \emph{flipped} collar $\bar F(x,y)=F(x,-y)$ is again a collar; everything below that is stated ``above'' the arc
holds ``below'' by flipping.

\formal{\lean{HPC.Hagopian.IsCollar}, \lean{HPC.Hagopian.exists_collar}, \lean{HPC.Hagopian.collarCoord},
\lean{HPC.Hagopian.flipC}.}

Fix $\kappa=1/32$. The \emph{middle band} is $X\in[\frac14+\kappa,\frac34-\kappa]=[\frac9{32},\frac{23}{32}]$; it contains the column
$[\frac38,\frac58]$. For $\eta\le1$ let $Q_\eta=F([\frac14,\frac34]\times[-\eta,\eta])$ be the \emph{box} of height $\eta$.

\begin{definition}\label{def:pass}
A \emph{pass} of $Q_\eta$ is a piece $\ell[a,b]$, $a<b$, with $\ell(a)$ on the left edge $F(\set{\frac14}\times[-\eta,\eta])$, $\ell(b)$ on the right
edge $F(\set{\frac34}\times[-\eta,\eta])$, and $\ell((a,b))$ in the open box $F((\frac14,\frac34)\times(-\eta,\eta))$. A pass has property
\emph{(Mon)} if $X(\ell(v))>X(\ell(v'))-2\kappa$ whenever $a\le v'<v\le b$. It is \emph{above} (below) the arc if it lies, end points
included, in $Y>0$ ($Y<0$).
\end{definition}

\formal{\lean{HPC.Hagopian.IsPass}, \lean{HPC.Hagopian.PassMon}, \lean{HPC.Hagopian.PassAbove},
\lean{HPC.Hagopian.PassBelow}.}

On a pass with (Mon), if $X(\ell v)=c$, $X(\ell v')=c'$ and $c'-c\ge2\kappa$, then $v<v'$ (the \emph{order corollary}).

\begin{lemma}[Passes]\label{lem:passes}
\leavevmode
\begin{enumerate}[label=(P\arabic*)]
\item Distinct passes of $Q_\eta$ have disjoint closed parameter intervals.
\item Only finitely many passes meet a bounded parameter interval; so among the passes with a given property starting after
a given time, if any, there is a first one.
\item $\ell[\frac14,\frac34]$ is a pass (the \emph{base pass}).
\item Every other pass is above or below the arc.
\end{enumerate}
For $\eta'\le\eta$, every pass of $Q_{\eta'}$ is a pass of $Q_\eta$.
\end{lemma}

\begin{proof}
(P1) If $a<a'<b$, then $\ell(a')$ would be on the left edge and in the open box; if $a'=b$, on both edges; if $a=a'$ and $b<b'$,
$\ell(b)$ would be on the right edge and in the open box. (P2) Passes inside $[-N,N]$ have length at least the modulus of uniform
continuity of $\ell$ for the distance between the edges, and at most two straddle $\pm N$. (P3) is clear. (P4) A pass other than the
base pass misses the base arc (injectivity of $\ell$ and (P1)), and its preimage under $F$ is connected.
\end{proof}

\formal{\lean{HPC.Hagopian.IsPass.disjoint}, \lean{HPC.Hagopian.IsPass.finite}, \lean{HPC.Hagopian.exists_first_pass},
\lean{HPC.Hagopian.isPass_base}, \lean{HPC.Hagopian.IsPass.above_or_below}, \lean{HPC.Hagopian.IsPass.mono}.}

\begin{lemma}[The shadow lemma]\label{lem:shadow}
For every $\eta'\in(0,1]$ there is $\delta>0$ such that: if $x\in[\frac9{32},\frac{23}{32}]$ and $d(\ell(u),\ell(x))<\delta$, then $u$ lies in the open
parameter interval of a pass of $Q_{\eta'}$ with (Mon).
\end{lemma}

\begin{proof}
Take $\varepsilon<\varepsilon_0$ (Theorem~\ref{thm:RO}) such that every point within $\varepsilon$ of $\gamma(s)$, $s\in[\kappa,1-\kappa]$, lies in the open collar with
$\abs{X-s}<\kappa$ and $\abs Y<\eta'$ (uniform continuity of $F^{-1}$ near the compact arc), and let $\delta$ be the Effros bound for $\varepsilon$.
\begin{enumerate}
\item The Effros theorem gives $h$ with $h(\ell(x))=\ell(u)$ moving points less than $\varepsilon$; by (O), applied to the shifted line, $h\circ\ell=\ell\circ\psi$
with $\psi$ increasing, continuous and onto, $\psi(x)=u$. Let $\theta=\psi^{-1}$.
\item Put $s_1=\frac3{16}$, $s_2=\frac{13}{16}$. For $v\in[\psi(s_1),\psi(s_2)]$ the point $\ell(v)=h(\gamma(\theta(v)))$ lies in the open collar with
$\abs{X(\ell v)-\theta(v)}<\kappa$ and $\abs Y<\eta'$. So $X<\frac14$ at $\psi(s_1)$ and $X>\frac34$ at $\psi(s_2)$.
\item Let $a=\max\set{v\in[\psi(s_1),u]:X(\ell v)\le\frac14}$ and $b=\min\set{v\in[u,\psi(s_2)]:X(\ell v)\ge\frac34}$. On $(a,u]$, $\frac14<X<\theta(v)+\kappa\le
x+\kappa\le\frac34$; on $[u,b)$, $\frac34>X>\theta(v)-\kappa\ge x-\kappa\ge\frac14$. Hence $a<u<b$, $X(\ell a)=\frac14$, $X(\ell b)=\frac34$, and
$\ell[a,b]$ is a pass of $Q_{\eta'}$.
\item (Mon): for $v'<v$, $X(\ell v)>\theta(v)-\kappa>\theta(v')-\kappa>X(\ell v')-2\kappa$.\qedhere
\end{enumerate}
\end{proof}

\formal{\lean{HPC.Hagopian.shadow_pass}.}

The shadow lemma holds only in the middle band: near an edge the line may turn back by up to $2\kappa$.

\begin{corollary}[Two scales]\label{cor:scales}
Put $\eta=\frac12$.
\begin{enumerate}[label=(\alph*)]
\item \emph{(Flow box.)} There is $\eta_2\in(0,\eta)$ such that every $u$ with $\ell(u)\in F([\frac9{32},\frac{23}{32}]\times[-\eta_2,\eta_2])$ lies in a pass of
$Q_\eta$ with (Mon); by (P1) every pass of $Q_\eta$ through such a point has (Mon).
\item A pass of $Q_\eta$ is \emph{thin} if it lies in $Q_{\eta_2}$. Every thin pass has (Mon) (it meets $X=\frac12$).
\item \emph{(Scale two.)} There is $\delta_2>0$ such that if $d(\ell(u),\ell(\frac12))<\delta_2$, then $u$ lies in the open interval of a thin pass.
\item \emph{(Late thin passes.)} For every $T$ there is a thin pass starting after $T$.
\end{enumerate}
\end{corollary}

\begin{proof}
(a) Take $\delta_1$ from Lemma~\ref{lem:shadow} for $\eta'=\eta$ and $\eta_2$ with $d(F(x,y),F(x,0))<\delta_1$ for $\abs y\le\eta_2$; apply the lemma at
$x=X(\ell u)$. (c) Lemma~\ref{lem:shadow} with $\eta'=\eta_2/2$, nesting and (P1). (d) Let $T'$ be the end of the pass containing $T$, if any,
and $T'=T$ otherwise. By recurrence (Theorem~\ref{thm:recurrent}) there is $u>T'$ with $d(\ell(u),\ell(\frac12))<\delta_2$; the thin pass $[a,b]\ni u$
has $a>T$, since otherwise it would contain $T$ and end at $T'<u$.
\end{proof}

\formal{\lean{HPC.Hagopian.flow_box}, \lean{HPC.Hagopian.FBProp}, \lean{HPC.Hagopian.PassThin},
\lean{HPC.Hagopian.scale_two}, \lean{HPC.Hagopian.exists_late_thin_pass}.}

\section{The trap}\label{sec:trap}

Let $P_m=\ell[a_m,b_m]$ be a thin pass above the arc with $a_m>\frac34$.

\begin{lemma}[The region under a thin pass]\label{lem:region}
$I_m=F^{-1}(P_m)$ is an arc from $r=(\frac14,r_2)$ to $s=(\frac34,s_2)$ in $[\frac14,\frac34]\times(0,\eta_2]$, meeting the vertical sides only in $r,s$.
There is a connected bounded open set $D_m\subseteq(\frac14,\frac34)\times(0,\eta_2)$ with
\[
\fr D_m=I_m\cup[\tfrac14,\tfrac34]\times\set0\cup\set{\tfrac14}\times[0,r_2]\cup\set{\tfrac34}\times[0,s_2].
\]
Put $R_m=F(D_m)$. For $c\in(\frac14,\frac34)$ let $w_m(c)$ be the least height of a point of $I_m$ over $c$ and $\tau_m(c)=\set c\times(0,w_m(c))$; let
$h_m>0$ be below all heights of $I_m$. Then $\tau_m(c)\subseteq D_m$, $(\frac14,\frac34)\times(0,h_m)\subseteq D_m$, and $\tau_m(c)$ is a cross-cut of $D_m$: no
connected subset of $D_m\setminus\tau_m(c)$ has closure meeting both open vertical sides.
\end{lemma}

\begin{proof}[Proof in outline]
$D_m$ is the bounded complementary domain of the simple closed curve formed by $I_m$ and the three sides, given by the Jordan curve
theorem (\vfive{Theorem~4.8}). For the cross-cut, the Jordan curve formed by the part of $I_m$ from $r$ to $(c,w_m(c))$, the left side, the
bottom from $\frac14$ to $c$ and $\tau_m(c)$ has the points of $D_m$ near the left side inside and the points near the right side outside.
\end{proof}

\formal{\lean{HPC.Hagopian.rect_domain_aux}, \lean{HPC.Hagopian.rect_domain_low}, \lean{HPC.Hagopian.seg_subset_rect_domain},
\lean{HPC.Hagopian.rect_domain_crosscut}, \lean{HPC.Hagopian.thinPass_arc}.}

\begin{lemma}[Passes entering the region]\label{lem:R1R2}
Let $P'=\ell[a',b']$ be a pass other than the base pass and $P_m$ whose interior meets $R_m$. Then (R1) its interior lies in $R_m$, its end
points lie on the open sides $F(\set{\frac14}\times(0,r_2))$ and $F(\set{\frac34}\times(0,s_2))$, and $P'$ is thin; and (R2) its interior meets $F(\tau_m(c))$
for every $c\in(\frac14,\frac34)$.
\end{lemma}

\begin{proof}
(R1) By (P4) $P'$ is above the arc; the preimage of its interior is connected, meets $D_m$, and misses $I_m$ (by (P1) and injectivity), the
bottom and the sides, so it lies in $D_m$. Its end points are limits of the interior, so they lie in $\fr D_m$, on the sides and not at the
corners. (R2) is the cross-cut property of Lemma~\ref{lem:region}.
\end{proof}

\formal{\lean{HPC.Hagopian.region_pass}, \lean{HPC.Hagopian.region_cross}.}

\begin{theorem}[The trapping lemma]\label{thm:trap}
Suppose that every thin pass above the arc starting after $\frac34$ starts at or after $a_m$. Then no pass starting after $b_m$ has its
interior in $R_m$.
\end{theorem}

\begin{proof}
\emph{Lemma A: no $v\in(\frac34,a_m)$ has $\ell(v)\in R_m$ with $X(\ell v)\in[\frac9{32},\frac{23}{32}]$.} Otherwise $R_m\subseteq Q_{\eta_2}$ and the flow box put $v$
in a pass $P'\ne$ base, $P_m$; by (R1) $P'$ is thin with interior in $R_m$, and by (P1) its interval lies in $(\frac34,a_m)$, against the
hypothesis.

\emph{Notation.} $w_L=w_m(\frac38)$, $y_m=w_m(\frac12)$, $w_R=w_m(\frac58)$; $T_L=F(\tau_m(\frac38))$, $\sigma=F(\set{\frac12}\times[0,y_m])$, $T_R=F(\tau_m(\frac58))$;
$x_m=F(\frac12,y_m)=\ell(t_m)$ and $z_R=F(\frac58,w_R)=\ell(t_R)$, with $t_m,t_R\in(a_m,b_m)$.

\emph{Step 1: $t_m<t_R$}, by (Mon) for $P_m$ and the order corollary ($\frac58-\frac12\ge2\kappa$).

\emph{Step 2: $J=\sigma\cup\ell[\frac12,t_m]$ is a simple closed curve.} $\sigma\cap\ell[\frac12,t_m]=\set{\ell(\frac12),x_m}$: the base piece $\ell[\frac12,\frac34]$
meets $\sigma$ only at $\ell(\frac12)$; the open segment $\sigma^\circ$ lies in $R_m$ and Lemma A excludes points of $\ell((\frac34,a_m))$ on it; points of $P_m$
on $X=\frac12$ have height at least $y_m$. By the Jordan curve theorem, $\R^2\setminus J=W\cup W'$ with $W,W'$ open, connected, disjoint, and
$\fr W=\fr W'=J$. Note \textbf{(J1)}: for $v>t_m$, $\ell(v)\in J$ iff $\ell(v)\in\sigma^\circ$.

\emph{Step 3.} Put $h=h_m$, $B_L=F([\frac38,\frac12)\times(0,h))$, $B_R=F((\frac12,\frac58]\times(0,h))$. The four connected sets $B_L,T_L,B_R,T_R$ lie in $R_m$
and miss $J$ (they miss $\sigma$ as $X\ne\frac12$, the base arc as $Y>0$, $P_m$ as they lie below its lowest points, and $\ell((\frac34,a_m))$ by
Lemma A). $B_L$ meets $T_L$ and $B_R$ meets $T_R$; let $W_L$ and $W_R$ be the components of $\R^2\setminus J$ containing $B_L\cup T_L$ and $B_R\cup T_R$.

\emph{Step 4: $W_L\ne W_R$.} Let $p=F(\frac12,h/2)\in\sigma^\circ$ and $S$ a small open square about $(\frac12,h/2)$ inside $[\frac38,\frac58]\times(0,h)$. $F(S)$ is
open and $F(S)\setminus J\subseteq B_L\cup B_R$. As $p$ is in the frontier of both components of $\R^2\setminus J$, each of them meets $B_L\cup B_R$; if both
$B_L$ and $B_R$ lay in one component, the other would miss $B_L\cup B_R$.

\emph{Step 5: $z_R\in W_R$.} $z_R\notin J$ by (J1), and $z_R\in\cl{T_R}$.

\emph{Step 6.} Suppose $P=\ell[a,b]$ has $a>b_m$ and interior in $R_m$. By (R2) at $c=\frac38$ there is $v_L\in(a,b)$ with $\ell(v_L)\in T_L\subseteq W_L$.
The path $\ell[t_R,v_L]$ joins $W_R$ to $W_L$, so it meets $J$; let $s^*$ be the largest $v\in[t_R,v_L]$ with $\ell(v)\in J$. By (J1),
$\ell(s^*)\in\sigma^\circ$, and $\ell((s^*,v_L])\subseteq W_L$. Also $s^*>b_m$, because $\ell((t_m,b_m])$ misses $\sigma$. By the flow box at $X=\frac12$, $s^*$
lies in the open interval of a pass $P^*=\ell[a^*,b^*]$ with (Mon); by (R1) its interior lies in $R_m$ and by (P1) $a^*>b_m$. If $s^*\in(a,v_L)$ then
$P^*=P$, and (Mon) for $P$ would put its point on $X=\frac38$ after its point on $X=\frac12$, which the order corollary forbids; so $b^*<a$. By
(R2) at $c=\frac58$ there is $v_R\in(a^*,b^*)$ with $\ell(v_R)\in T_R\subseteq W_R$, and (Mon) for $P^*$ gives $v_R>s^*$. So $v_R\in(s^*,v_L)$ and
$\ell(v_R)\in W_L$, contradicting Step~4.
\end{proof}

\formal{\lean{HPC.Hagopian.region_lemmaA}, \lean{HPC.Hagopian.two_sides}, \lean{HPC.Hagopian.trap_jordan},
\lean{HPC.Hagopian.trap_pass}; the Jordan curve theorem is \lean{HPC.Jones.jordan_curve_theorem}.}

\section{The contradiction}

\begin{theorem}[Bing 1960; Jones 1975]\label{thm:noarc}
No homogeneous indecomposable plane continuum contains an arc.
\end{theorem}

\begin{proof}
Suppose $M$ does, and let $\ell$, $F$ be as above.

\emph{C1.} By Corollary~\ref{cor:scales}(d) and (P2) let $P_1=\ell[a_1,b_1]$ be the first thin pass with $a_1>\frac34$; flipping the collar if
necessary, it is above the arc. The hypothesis of Theorem~\ref{thm:trap} holds for $R_1$: a violating pass would be thin (R1) and start in
$(\frac34,a_1)$.

\emph{C2.} Take $h_1=h_{m_1}$ and $\rho_1$ with $B(\ell(\frac12),\rho_1)\subseteq F((\frac38,\frac58)\times(-h_1,h_1))$. Recurrence gives $u_1>b_1$ with
$d(\ell(u_1),\ell(\frac12))<\min(\delta_2,\rho_1)$.

\emph{C3.} Write $\ell(u_1)=F(x,y)$. If $y=0$, then $u_1=x<b_1$ by injectivity. If $y>0$, $\ell(u_1)\in R_1$ lies on a thin pass (scale two) whose
interior is in $R_1$ (R1) and which starts after $b_1$ (P1), against Theorem~\ref{thm:trap}. So $y<0$, and some thin pass below the arc starts after
$b_1$; let $P_2=\ell[a_2,b_2]$ be the first.

\emph{C4.} The trapping lemma applies to $P_2$ (with the flipped collar): a thin pass below the arc starting in $(\frac34,a_2)$ would, by (P1),
start in $(\frac34,a_1)$, against the minimality of $P_1$, or in $(b_1,a_2)$, against the minimality of $P_2$.

\emph{C5.} Take $h=\min(h_1,h_2)$ and $\rho$ with $B(\ell(\frac12),\rho)\subseteq F((\frac38,\frac58)\times(-h,h))$; recurrence gives $u>b_2$ with
$d(\ell(u),\ell(\frac12))<\min(\delta_2,\rho)$. $Y(\ell u)=0$ is excluded by injectivity; $Y>0$ puts $\ell(u)$ in $R_1$ on a thin pass after $b_1$ (C1);
$Y<0$ puts it in $R_2$ on a thin pass after $b_2$ (C4). This is a contradiction.
\end{proof}

\formal{\lean{HPC.Hagopian.jones75_final_aux}, \lean{HPC.Hagopian.jones75_contra},
\lean{HPC.Hagopian.NoArcIndecHomPlane}, \lean{HPC.Hagopian.jones75_noArc}.}

\section{Hereditary indecomposability}

\begin{lemma}\label{lem:monotone-image}
A monotone image of an indecomposable continuum is indecomposable: if $q\colon X\to Y$ is a continuous surjection of a continuum onto a metric space
with connected fibres and $X$ is indecomposable, so is $Y$.
\end{lemma}

\begin{proof}
$q$ is closed with connected fibres, so preimages of subcontinua are subcontinua (as in Lemma~\ref{lem:adm-basic}(a)). If $Y=H\cup K$ with $H,K$
proper subcontinua, then $X=q^{-1}(H)\cup q^{-1}(K)$ with both proper.
\end{proof}

\formal{\lean{HPC.Hagopian.indecomposable_of_monotone_image}.}

\begin{theorem}[Hagopian 1976, Theorem~1]\label{thm:HI}
Every homogeneous indecomposable plane continuum is hereditarily indecomposable.
\end{theorem}

\begin{proof}
This is Hagopian's 1984 proof (his Corollary~3), following a plan of Jones. Suppose $M$ is not hereditarily indecomposable, and let $q\colon M\to Y$ be
as in Theorem~\ref{thm:h84}. Its fibres are proper subcontinua of $M$, so they do not separate the plane (Theorem~\ref{thm:nosep}). By Moore's
theorem on monotone decompositions (\vfive{Corollary~9.16}), $Y$ is homeomorphic to a compact plane set $Q$. $Q$ is connected (an image of $M$),
homogeneous (as $Y$ is), indecomposable (Lemma~\ref{lem:monotone-image}) and contains an arc. This contradicts Theorem~\ref{thm:noarc}.
\end{proof}

\formal{\lean{HPC.Hagopian.exists_plane_quotient_of_monotone}, \lean{HPC.Hagopian.hagopian76_theorem1_of_jones},
\lean{HPC.Hagopian.hagopian76_theorem1_jones}, \lean{HPC.Hagopian.Hagopian76Thm1}, \lean{HPC.Hagopian.hagopian76_theorem1}.}

\begin{corollary}\label{cor:bing60}
\leavevmode
\begin{enumerate}[label=(\alph*)]
\item \emph{(Bing 1960, for plane continua.)} A homogeneous plane continuum that contains an arc is a simple closed curve.
\item \emph{(Hagopian 1976, Theorem~2.)} A homogeneous plane continuum that does not separate the plane is hereditarily indecomposable.
\end{enumerate}
\end{corollary}

These consequences are recorded in the formal development; they are not used in this volume.

\formal{\lean{HPC.Hagopian.bing60_plane_uncond}, \lean{HPC.Hagopian.hagopian76_theorem2_uncond}.}

\section*{Exercises}

\begin{exercise}
Show that in a compact metric space the forward limit set of a continuous map $\ell\colon\R\to Y$ is nonempty and compact, and that it does not
change if $\ell$ is replaced by $s\mapsto\ell(s+c)$.
\end{exercise}

\begin{exercise}
Show that the dyadic solenoid is homogeneous and that its arc components are lines. Which part of the proof of Theorem~\ref{thm:noarc} fails for
it?
\end{exercise}
